

\dresumo{
A criação de pipelines de dados por meio de soluções proprietárias apresenta algumas desvantagens competitivas, como a falta de transparência, a dependência de fornecedores e, principalmente, os altos custos. O \textit{software} de código aberto surge nesse contexto com um modelo de desenvolvimento descentralizado e colaborativo, que distribui o código-fonte publicamente, permitindo que qualquer pessoa possa usar, alterar e redistribuir como preferir, sem custo adicional. No entanto, o mercado ainda enfrenta desafios e resistências na adoção dessas soluções. Para enfrentar esse problema, este trabalho mapeou os principais \textit{softwares} livres utilizados no mercado para atividades de engenharia de dados, com o propósito de criar um catálogo de aplicações para construção de \textit{pipeline} de dados \textit{open source}. Com o intuito de demonstrar a viabilidade prática, foi desenvolvido um exemplo prático de \textit{pipeline} com algumas das ferramentas encontradas, desde a concepção inicial até a implementação final, aplicado a um cenário de coleta de informações sobre passagens aéreas para o Brasil, extraindo as informações de diferentes sites internacionais. A análise focou no impacto dessas ferramentas em termos de escalabilidade e desempenho, com o intuito de oferecer uma base técnica relevante para trabalhos acadêmicos e pequenas e médias empresas que buscam soluções eficientes e acessíveis, sem comprometer seus recursos.
}
{\textit{Pipeline} de Dados; \textit{Software} Livre; Engenharia de Dados}

\dabstract{
    The development of data pipelines utilising proprietary solutions entails several competitive drawbacks, including opacity, vendor dependency, and, crucially, elevated expenses. Open-source software arises inside this framework, characterised by a decentralised and collaborative development process that publicly disseminates the source code, permitting anybody to utilise, alter, and redistribute it at no additional cost. Nonetheless, the market continues to encounter obstacles and opposition in embracing these solutions. This study catalogued the primary free software utilised in the market for data engineering tasks, aiming to establish a repository of apps for constructing open-source data pipelines. To illustrate practical viability, a pipeline was constructed utilising various tools, from initial conception to final implementation, focused on gathering data regarding airline tickets to Brazil by extracting information from multiple international websites. The investigation examined the influence of these tools regarding scalability and performance, intending to furnish a pertinent technical foundation for academic research and small to medium-sized firms in pursuit of efficient and cost-effective solutions without depleting their resources.
}
{\textit{Data Pipeline}; \textit{Open Source}; \textit{Data Engineering}}